% Part: first-order-logic
% Chapter: natural-deduction
% Section: soundness

\documentclass[../../../include/open-logic-section]{subfiles}

\begin{document}

\iftag{FOL}
      {\olfileid{fol}{ntd}{sou}}
      {\olfileid{pl}{ntd}{sou}}
      
\olsection{可靠性}

\begin{explain}
一个推导系统（如自然演绎），如果它不能推出并非实际后承的结论，就是\emph{可靠的}。因此，可靠性是推导系统的一种保证其安全性的性质。视所讨论的证明论性质而定，我们希望知道例如：
\begin{enumerate}
\item 每个可推导的语句都是\iftag{FOL}{有效的}{重言式}；
\item 如果一个语句可以从另一些语句推导出来，那么它也是这些语句的后承；
\item 如果一个语句集合不一致，那么它是不可满足的。
\end{enumerate}
这些是推导系统的重要性质。如果其中任何一个不成立，这个推导系统就是有缺陷的——它会推出太多东西。因此，确立一个推导系统的可靠性至关重要。
\end{explain}

\begin{thm}[可靠性]
\ollabel{thm:soundness}
如果 $!A$ 可由未解除的假设 $\Gamma$ 推导出来，那么 $\Gamma$ 语义地后承 $!A$。
\end{thm}

\begin{proof}
设 $\delta$ 是 $!A$ 的一个推导。我们对~$\delta$ 中的推理数目进行归纳。

对于归纳基础，我们证明当推理数目为~$0$ 时结论成立。此时，$\delta$ 仅由一个语句~$!A$ 构成，即一个假设。因为假设只能通过推理来解除，而这里没有推理，所以该假设是未解除的。因此，任何满足该证明所有未解除假设的\iftag{FOL}{结构~$\Struct{M}$}{赋值~$\pAssign{v}$}也满足~$!A$。

现在考虑归纳步骤。假设 $\delta$ 包含~$n$ 个推理。最末推理的前提是通过子推导导出的，每个子推导包含的推理都少于~$n$ 个。我们假定归纳假设成立：最末推理的前提是以这些前提为结论的子推导的未解除假设的后承。我们需要证明，结论~$!A$ 是整个证明的未解除假设的后承。

我们根据最末推理的类型来区分情况。首先考虑只有一个前提的可能推理。

\begin{enumerate}
\item 假设最后一个推理是 \Intro{\lnot}：推导形如
  \begin{prooftree}
    \AxiomC{$\Gamma, \Discharge{!A}{n}$}
    \RightLabel{$\delta_1$}
    \DeduceC{$\lfalse$}
    \DischargeRule{\Intro{\lnot}}{n}
    \UnaryInfC{$\lnot !A$}
  \end{prooftree}
  根据归纳假设，$\lfalse$ 是~$\delta_1$ 的未解除假设 $\Gamma \cup \{!A\}$ 的后承。考虑\iftag{FOL}{结构~$\Struct{M}$}{赋值~$\pAssign{v}$}。我们需要证明，如果 $\iftag{FOL}{\Sat{M}}{\pSat{v}}{\Gamma}$，则 $\iftag{FOL}{\Sat{M}}{\pSat{v}}{\lnot !A}$。用归谬法，假设 $\iftag{FOL}{\Sat{M}}{\pSat{v}}{\Gamma}$，但 $\iftag{FOL}{\Sat/{M}}{\pSat/{v}}{\lnot !A}$，即 $\iftag{FOL}{\Sat{M}}{\pSat{v}}{!A}$。这意味着 $\iftag{FOL}{\Sat{M}}{\pSat{v}}{\Gamma \cup \{!A\}}$。这与我们的归纳假设矛盾。所以，$\iftag{FOL}{\Sat{M}}{\pSat{v}}{\lnot !A}$。

\item 最后一个推理是 \Elim{\land}：有两种变体：可以从前提 $!A \land !B$ 推出 $!A$ 或 $!B$。考虑第一种情况。推导~$\delta$ 形如：
  \begin{prooftree}
    \AxiomC{$\Gamma$}
    \RightLabel{$\delta_1$}
    \DeduceC{$!A \land !B$}
    \RightLabel{\Elim{\land}}
    \UnaryInfC{$!A$}
  \end{prooftree}
  根据归纳假设，$!A \land !B$ 是~$\delta_1$ 的未解除假设~$\Gamma$ 的后承。考虑\iftag{FOL}{结构~$\Struct{M}$}{赋值~$\pAssign{v}$}。我们需要证明，如果 $\iftag{FOL}{\Sat{M}}{\pSat{v}}{\Gamma}$，则 $\iftag{FOL}{\Sat{M}}{\pSat{v}}{!A}$。假设 $\iftag{FOL}{\Sat{M}}{\pSat{v}}{\Gamma}$。根据归纳假设（$\Gamma \Entails !A \land !B$），我们知道 $\iftag{FOL}{\Sat{M}}{\pSat{v}}{!A \land !B}$。根据定义，$\iftag{FOL}{\Sat{M}}{\pSat{v}}{!A \land !B}$ 当且仅当 $\iftag{FOL}{\Sat{M}}{\pSat{v}}{!A}$ 且 $\iftag{FOL}{\Sat{M}}{\pSat{v}}{!B}$。（从 $!A \land !B$ 推出 $!B$ 的情况类似处理。）
  
\item 最后一个推理是 \Intro{\lor}：有两种变体：$!A \lor !B$ 可以从前提~$!A$ 或前提~$!B$ 推出。考虑第一种情况。推导形如：
  \begin{prooftree}
    \AxiomC{$\Gamma$}
    \RightLabel{$\delta_1$}
    \DeduceC{$!A$}
    \RightLabel{\Intro{\lor}}
    \UnaryInfC{$!A \lor !B$}
  \end{prooftree}
  根据归纳假设，$!A$ 是~$\delta_1$ 的未解除假设~$\Gamma$ 的后承。考虑\iftag{FOL}{结构~$\Struct{M}$}{赋值~$\pAssign{v}$}。我们需要证明，如果 $\iftag{FOL}{\Sat{M}}{\pSat{v}}{\Gamma}$，则 $\iftag{FOL}{\Sat{M}}{\pSat{v}}{!A \lor !B}$。假设 $\iftag{FOL}{\Sat{M}}{\pSat{v}}{\Gamma}$；由于 $\Gamma \Entails !A$（归纳假设），则 $\iftag{FOL}{\Sat{M}}{\pSat{v}}{!A}$。因此也必有 $\iftag{FOL}{\Sat{M}}{\pSat{v}}{!A \lor !B}$。（从~$!B$ 推出 $!A \lor !B$ 的情况类似处理。）
  
\item 最后一个推理是 \Intro{\lif}：$!A \lif !B$ 是从一个以~$!A$ 为假设、以~$!B$ 为结论的子推导推出的，即：
  \begin{prooftree}
    \AxiomC{$\Gamma, \Discharge{!A}{n}$}
    \RightLabel{$\delta_1$}
    \DeduceC{$!B$}
    \DischargeRule{\Intro{\lif}}{n}
    \UnaryInfC{$!A \lif !B$}
  \end{prooftree}
  根据归纳假设，$!B$ 是~$\delta_1$ 的未解除假设的后承，即 $\Gamma \cup \{!A\} \Entails
  !B$。考虑\iftag{FOL}{结构~$\Struct{M}$}{赋值~$\pAssign{v}$}。$\delta$ 的未解除假设仅为 $\Gamma$，因为 $!A$ 在最后一个推理中被解除了。所以我们需要证明 $\Gamma \Entails !A \lif !B$。用归谬法，假设对某个\iftag{FOL}{结构~$\Struct{M}$}{赋值~$\pAssign{v}$}，$\iftag{FOL}{\Sat{M}}{\pSat{v}}{\Gamma}$ 但 $\iftag{FOL}{\Sat/{M}}{\pSat/{v}}{!A \lif !B}$。所以，$\iftag{FOL}{\Sat{M}}{\pSat{v}}{!A}$ 且 $\iftag{FOL}{\Sat/{M}}{\pSat/{v}}{!B}$。但根据归纳假设，$!B$ 是 $\Gamma \cup \{!A\}$ 的后承，即 $\iftag{FOL}{\Sat{M}}{\pSat{v}}{!B}$，矛盾。因此，$\Gamma \Entails !A \lif !B$。

\item 最后一个推理是 \FalseInt：这里，$\delta$ 以如下方式结束：
  \begin{prooftree}
    \AxiomC{$\Gamma$}
    \RightLabel{$\delta_1$}
    \DeduceC{$\lfalse$}
    \RightLabel{\FalseInt}
    \UnaryInfC{$!A$}
  \end{prooftree}
  根据归纳假设，$\Gamma \Entails \lfalse$。我们需要证明 $\Gamma \Entails !A$。假设不然；则对某个~$\iftag{FOL}{\Struct{M}}{\pAssign{v}}$ 有 $\iftag{FOL}{\Sat{M}}{\pSat{v}}{\Gamma}$ 和 $\iftag{FOL}{\Sat/{M}}{\pSat/{v}}{!A}$。但我们总有 $\iftag{FOL}{\Sat/{M}}{\pSat/{v}}{\lfalse}$，所以这意味着 $\Gamma \Entails/ \lfalse$，与归纳假设矛盾。

\item 最后一个推理是 \FalseCl：练习。

\iftag{FOL}{
\item 最后一个推理是 \Intro{\lforall}：那么 $\delta$ 形如
  \begin{prooftree}
    \AxiomC{$\Gamma$}
    \RightLabel{$\delta_1$}
    \DeduceC{$!A(a)$}
    \RightLabel{\Intro{\lforall}}
    \UnaryInfC{$\lforall[x][!A(x)]$}
  \end{prooftree}
  根据归纳假设，前提 $!A(a)$ 是未解除假设 $\Gamma$ 的后承。考虑某个满足 $\Sat{M}{\Gamma}$ 的结构 $\Struct{M}$。我们需要证明 $\Sat{M}{\lforall[x][!A(x)]}$。由于 $\lforall[x][!A(x)]$ 是语句，这意味着我们必须证明对每个变元指派~$s$，$\Sat{M}{!A(x)}[s]$ (\olref[syn][ass]{prop:sat-quant})。因为 $\Gamma$ 完全由语句组成，根据 \olref[syn][sat]{defn:satisfaction}，对任意 $!B \in \Gamma$，$\Sat{M}{!B}[s]$。令 $\Struct{M'}$ 除了 $\Assign{a}{M'} = s(x)$ 外与 $\Struct{M}$ 相同。由于 $a$ 不在~$\Gamma$ 中出现，根据 \olref[syn][ext]{cor:extensionality-sent}，有 $\Sat{M'}{\Gamma}$。因为 $\Gamma \Entails !A(a)$，所以 $\Sat{M'}{!A(a)}$。由于 $!A(a)$ 是语句，根据 \olref[syn][ass]{prop:sentence-sat-true}，$\Sat{M'}{!A(a)}[s]$。根据 \olref[syn][ext]{prop:ext-formulas}，$\Sat{M'}{!A(x)}[s]$ 当且仅当 $\Sat{M'}{!A(a)}$（回想 $!A(a)$ 就是 $\Subst{!A(x)}{a}{x}$）。所以，$\Sat{M'}{!A(x)}[s]$。由于 $a$ 不在 $!A(x)$ 中出现，根据 \olref[syn][ext]{prop:extensionality}，$\Sat{M}{!A(x)}[s]$。但 $s$ 是任意变元指派，所以 $\Sat{M}{\lforall[x][!A(x)]}$。
  
\item 最后一个推理是 \Intro{\lexists}：练习。

\item 最后一个推理是 \Elim{\forall}：练习。
}{}
\end{enumerate}

现在考虑有多个前提的可能推理：\Elim{\lor}、\Intro{\land}、\iftag{FOL}{\Elim{\lif} 和 \Elim{\lexists}}{和 \Elim{\lif}}。
\begin{enumerate}

\item 最后一个推理是 \Intro{\land}。$!A \land !B$ 是从前提 $!A$ 和 $!B$ 推出的，且 $\delta$ 形如：
  \begin{prooftree}
    \AxiomC{$\Gamma_1$}
    \RightLabel{$\delta_1$}
    \DeduceC{$!A$}
    \AxiomC{$\Gamma_2$}
    \RightLabel{$\delta_2$}
    \DeduceC{$!B$}
    \RightLabel{\Intro{\land}}
    \BinaryInfC{$!A \land !B$}
  \end{prooftree}
  根据归纳假设，$!A$ 是~$\delta_1$ 的未解除假设~$\Gamma_1$ 的后承，且 $!B$ 是~$\delta_2$ 的未解除假设~$\Gamma_2$ 的后承。$\delta$ 的未解除假设是 $\Gamma_1 \cup \Gamma_2$，所以我们需要证明 $\Gamma_1 \cup \Gamma_2 \Entails
  !A \land !B$。考虑满足 $\iftag{FOL}{\Sat{M}}{\pSat{v}}{\Gamma_1 \cup \Gamma_2}$ 的\iftag{FOL}{结构~$\Struct{M}$}{赋值~$\pAssign{v}$}。由于 $\iftag{FOL}{\Sat{M}}{\pSat{v}}{\Gamma_1}$，且 $\Gamma_1 \Entails !A$，所以必有 $\iftag{FOL}{\Sat{M}}{\pSat{v}}{!A}$；同时由于 $\iftag{FOL}{\Sat{M}}{\pSat{v}}{\Gamma_2}$，且 $\Gamma_2 \Entails !B$，所以有 $\iftag{FOL}{\Sat{M}}{\pSat{v}}{!B}$。合起来，$\iftag{FOL}{\Sat{M}}{\pSat{v}}{!A \land !B}$。
  
\item 最后一个推理是 \Elim{\lor}：练习。

\item 最后一个推理是 \Elim{\lif}。$!B$ 是从前提 $!A \lif !B$ 和~$!A$ 推出的。推导~$\delta$ 形如：
  \begin{prooftree}
    \AxiomC{$\Gamma_1$}
    \RightLabel{$\delta_1$}
    \DeduceC{$!A \lif !B$}
    \AxiomC{$\Gamma_2$}
    \RightLabel{$\delta_2$}
    \DeduceC{$!A$}
    \RightLabel{\Elim{\lif}}
    \BinaryInfC{$!B$}
  \end{prooftree}
  根据归纳假设，$!A \lif !B$ 是~$\delta_1$ 的未解除假设~$\Gamma_1$ 的后承，且 $!A$ 是~$\delta_2$ 的未解除假设~$\Gamma_2$ 的后承。考虑\iftag{FOL}{结构~$\Struct{M}$}{赋值~$\pAssign{v}$}。我们需要证明，如果 $\iftag{FOL}{\Sat{M}}{\pSat{v}}{\Gamma_1 \cup
    \Gamma_2}$，则 $\iftag{FOL}{\Sat{M}}{\pSat{v}}{!B}$。假设 $\iftag{FOL}{\Sat{M}}{\pSat{v}}{\Gamma_1 \cup \Gamma_2}$。由于 $\Gamma_1 \Entails !A \lif !B$，所以 $\iftag{FOL}{\Sat{M}}{\pSat{v}}{!A
    \lif !B}$。由于 $\Gamma_2 \Entails !A$，我们有 $\iftag{FOL}{\Sat{M}}{\pSat{v}}{!A}$。这意味着 $\iftag{FOL}{\Sat{M}}{\pSat{v}}{!B}$（因为如果 $\iftag{FOL}{\Sat/{M}}{\pSat/{v}}{!B}$，由于 $\iftag{FOL}{\Sat{M}}{\pSat{v}}{!A}$，我们将有 $\iftag{FOL}{\Sat/{M}}{\pSat/{v}}{!A \lif !B}$，这与~$\iftag{FOL}{\Sat{M}}{\pSat{v}}{!A \lif !B}$矛盾）。

\item 最后一个推理是 \Elim{\lnot}：练习。
  
\tagitem{FOL}{最后一个推理是 \Elim{\lexists}：练习。}{}
\end{enumerate}
\end{proof}

\tagprob{FOL}


\begin{prob}
完成 \olref[fol][ntd][sou]{thm:soundness} 的证明。
\end{prob}


\tagendprob

\tagprob{notFOL}


\begin{prob}
完成 \olref[pl][ntd][sou]{thm:soundness} 的证明。
\end{prob}


\tagendprob

\begin{cor}
\ollabel{cor:weak-soundness}
若 $\Proves !A$，则 $!A$ 是 \iftag{FOL}{有效的}{重言式}。
\end{cor}

\begin{cor}
\ollabel{cor:consistency-soundness}
若 $\Gamma$ 可满足，则 $\Gamma$ 是一致的。
\end{cor}

\begin{proof}
我们证明其逆否命题。
假设 $\Gamma$ 不一致。
则 $\Gamma \Proves \lfalse$，即存在一个从 $\Gamma$ 中的未解除假设推导出 $\lfalse$ 的推导。
根据 \olref{thm:soundness}，任何满足 $\Gamma$ 的
\iftag{FOL}{结构~$\Struct{M}$}{赋值~$\pAssign{v}$}
也必须满足 $\lfalse$。
由于 $\iftag{FOL}{\Sat/{M}}{\pSat/{v}}{\lfalse}$ 对每个
\iftag{FOL}{结构~$\Struct{M}$}{赋值~$\pAssign{v}$} 都成立，
所以没有 \iftag{FOL}{$\Struct{M}$}{$\pAssign{v}$} 能满足 $\Gamma$，
即 $\Gamma$ 不可满足。
\end{proof}